\GalSourcePageBoundary{076}{L0075}{47}{47}

Maintenant, le groupe qui précédera immédiatement celui-ci
dans l'ordre des décompositions devra se composer d'un certain
nombre de groupes ayant tous les mêmes substitutions que
celui-ci. Or, j'observe que ces substitutions peuvent s'exprimer
ainsi (faisons, en général, $x_{n}=x_{0}$, $x_{n+1}=x_{1}$, $\ldots$; il est clair
que chacune des substitutions du groupe (G) s'obtient en mettant
partout à la place de $x_{k}$, $x_{k+c}$, $c$ étant une constante).

Considérons l'un quelconque des groupes semblables au
groupe (G). D'après le théorème II, il devra s'obtenir en opérant
partout dans ce groupe une même substitution; par exemple, en
mettant partout dans le groupe (G), à la place de $x_{k}$, $x_{f(k)}$, $f$ étant
une certaine fonction.

Les substitutions de ces nouveaux groupes devant être les
mêmes que celles du groupe (G), on devra avoir
\[
\GalControlReading{f(k+c)=f(k)+C}{W09-CD003},
\]
$C$ étant indépendant de $k$.

Donc
\[
\begin{aligned}
f(k+2c)&=f(k)+2C,\\
&\hspace{2em}\dotfill\\
f(k+mc)&=f(k)+mC.
\end{aligned}
\]
Si $c=1$, $k=0$, on trouvera
\[
f(m)=am+b
\]
ou bien
\[
f(k)=ak+b,
\]
$a$ et $b$ étant des constantes.

Donc, le groupe qui précède immédiatement le groupe (G) ne
devra contenir que des substitutions telles que
\[
x_{k},\qquad x_{ak+b},
\]
et ne contiendra pas, par conséquent, d'autre substitution circulaire
que celle du groupe (G).
